\documentclass[11pt,a4paper]{article}

% ── Codificación, fuentes y lenguaje ───────────────────────────────────────────
\usepackage[utf8]{inputenc}
\usepackage[T1]{fontenc}
\usepackage[default,scale=0.95]{sourcesanspro}
\renewcommand{\familydefault}{\sfdefault}
\usepackage[spanish,es-noshorthands,es-tabla]{babel}

% ── Geometría y espaciado ─────────────────────────────────────────────────────
\usepackage[top=2.2cm, bottom=2.5cm, left=2.2cm, right=2.2cm, headheight=15pt]{geometry}
\usepackage{setspace}
\setstretch{1.18}
\usepackage{parskip}

% ── Matemáticas y unidades ────────────────────────────────────────────────────
\usepackage{amsmath}
\usepackage{amssymb}
\usepackage{siunitx}

% ── Circuitos ─────────────────────────────────────────────────────────────────
\usepackage[european, straightvoltages]{circuitikz}

% ── Tablas ────────────────────────────────────────────────────────────────────
\usepackage{booktabs}
\usepackage{tabularx}
\usepackage{array}
\usepackage{longtable}
\usepackage{colortbl}
\usepackage{enumitem}

% ── Colores, cajas e iconos ───────────────────────────────────────────────────
\usepackage[dvipsnames]{xcolor}
\usepackage{tcolorbox}
\tcbuselibrary{skins, breakable}
\usepackage{fontawesome5}
\usepackage{graphicx}

% ── Listados de código (dark-mode infográfico) ────────────────────────────────
\usepackage{listings}

% ── Secciones con barra de color (infografía) ────────────────────────────────
\usepackage{titlesec}
\titleformat{\section}
  {\normalfont\LARGE\bfseries\color{azulNoche}}
  {\colorbox{cyanNeon}{\color{fondoOscuro}\thesection}}{0.7em}{}
  [\vspace{2pt}\color{cyanNeon}\rule{\linewidth}{1.8pt}]
\titlespacing*{\section}{0pt}{24pt}{10pt}
\titleformat{\subsection}
  {\normalfont\large\bfseries\color{azulNoche}}
  {\textcolor{cyanNeon}{\thesubsection}}{0.7em}{}
  [\vspace{1pt}\color{grisLinea}\rule{\linewidth}{0.6pt}]
\titlespacing*{\subsection}{0pt}{14pt}{6pt}
\titleformat{\subsubsection}
  {\normalfont\normalsize\bfseries\color{azulNoche}}
  {\textcolor{cyanNeon}{\thesubsubsection}}{0.7em}{}
\titlespacing*{\subsubsection}{0pt}{10pt}{4pt}

% ── Cabeceras y pies de página ────────────────────────────────────────────────
\usepackage{fancyhdr}
\usepackage{lastpage}
\pagestyle{fancy}
\fancyhf{}
\renewcommand{\headrulewidth}{0pt}
\renewcommand{\footrulewidth}{0pt}
\fancyfoot[C]{%
  \begin{minipage}{\linewidth}
    \vspace{2pt}
    \color{cyanNeon}\rule{\linewidth}{1.2pt}\\[2pt]
    \centering\color{grisTexto}\footnotesize
    Página \thepage\ de \pageref{LastPage}
  \end{minipage}%
}

% ── Hiperenlaces ──────────────────────────────────────────────────────────────
\usepackage[colorlinks=true, linkcolor=azulNoche, urlcolor=cyanNeon]{hyperref}
\sloppy
\emergencystretch 3em

% ── Paleta de colores — tecnológico dark-mode ──────────────────────────────────
\definecolor{fondoOscuro}{HTML}{0D1B2A}     % Dark navy — fondo banda título
\definecolor{verdeTurquesa}{HTML}{004D40}    % Verde turquesa — bandas de marca
\definecolor{azulNoche}{HTML}{1B3A6B}       % Midnight blue — secciones, marcos
\definecolor{azulMedio}{HTML}{2A5298}       % Azul medio — variante de acento
\definecolor{cyanNeon}{HTML}{00C8E0}        % Cyan eléctrico — acentos primarios
\definecolor{verdeSignal}{HTML}{00B887}     % Verde señal — fortalezas, éxito
\definecolor{naranjaVivo}{HTML}{FF7043}     % Naranja vivo — mejoras, advertencias
\definecolor{rojoAlerta}{HTML}{E53935}      % Rojo alerta — errores
\definecolor{grisPapel}{HTML}{F4F6FA}       % Fondo tarjetas claras
\definecolor{grisLinea}{HTML}{C8D0E0}       % Bordes sutiles
\definecolor{grisTexto}{HTML}{6B7A99}       % Texto secundario
\definecolor{amarilloNota}{HTML}{FFB300}    % Notas / advertencias doradas

% Colores específicos para bloques de código (dark-mode)
\definecolor{codigoFondo}{HTML}{1E2D3D}      % Fondo del bloque de código
\definecolor{codigoComentario}{HTML}{7A8FA6} % Comentarios
\definecolor{codigoCadena}{HTML}{FF9D5C}     % Cadenas / strings
\definecolor{codigoKeyword}{HTML}{00C8E0}    % Palabras clave
\definecolor{codigoNumero}{HTML}{00E5A0}     % Números

% ── Banda de título: portada infográfica ──────────────────────────────────────
%   Diseño en 2 capas: banda de color (por defecto fondo oscuro) + franja de
%   acento cyan abajo. Uso: \bandaTitulo[color]{título}{subtítulo}.
%   El icono se puede cambiar con \renewcommand{\iconoBanda}{...} (FontAwesome).
\newcommand{\iconoBanda}{microchip}
\newcommand{\bandaTitulo}[3][fondoOscuro]{%
  \begin{tcolorbox}[
    enhanced,
    colback=#1, colframe=#1,
    boxrule=0pt, arc=8pt, outer arc=8pt,
    width=\linewidth,
    top=18pt, bottom=6pt, left=14pt, right=14pt,
    borderline south={3.5pt}{0pt}{cyanNeon},
  ]
    \begin{minipage}[c]{0.07\linewidth}
      \centering
      \textcolor{cyanNeon}{\fontsize{30}{30}\selectfont\faIcon{\iconoBanda}}
    \end{minipage}%
    \hspace{8pt}%
    \begin{minipage}[c]{0.88\linewidth}
      {\fontsize{20}{24}\selectfont\bfseries\color{white}#2}\par\vspace{5pt}%
      \textcolor{cyanNeon!80!white}{\small\faIcon{angle-right}~#3}%
    \end{minipage}
    \vspace{4pt}
  \end{tcolorbox}%
  \vspace{8pt}%
}

% ── Tarjetas de datos (infografía) ────────────────────────────────────────────
\tcbset{
  tarjetaDato/.style={
    enhanced, breakable,
    arc=6pt, outer arc=6pt,
    colback=grisPapel, colframe=azulNoche,
    boxrule=1pt,
    fonttitle=\bfseries\small, coltitle=white,
    attach boxed title to top left={yshift=-3mm, xshift=6mm},
    boxed title style={
      colback=azulNoche, colframe=azulNoche,
      arc=4pt, boxrule=0pt,
    },
    drop fuzzy shadow=azulNoche!30!white,
    top=8pt, bottom=8pt, left=10pt, right=10pt,
  },
  tarjetaIcono/.style={
    enhanced, breakable,
    arc=6pt, outer arc=6pt,
    colback=white, colframe=grisLinea, boxrule=0.8pt,
    fonttitle=\bfseries\small, coltitle=azulNoche,
    borderline west={3pt}{0pt}{cyanNeon},
    drop fuzzy shadow=azulNoche!25!white,
    top=6pt, bottom=6pt, left=10pt, right=10pt,
  },
  cajaTitulo/.style={
    enhanced, breakable,
    arc=6pt, outer arc=6pt,
    colback=azulNoche!10!grisPapel, colframe=azulNoche,
    boxrule=1pt,
    fonttitle=\bfseries, coltitle=white,
    attach boxed title to top left={yshift=-3mm, xshift=6mm},
    boxed title style={
      colback=azulNoche, colframe=azulNoche,
      arc=4pt, boxrule=0pt,
    },
    drop fuzzy shadow=azulNoche!25!white,
    top=8pt, bottom=8pt, left=10pt, right=10pt,
  },
  cajaContenido/.style={
    enhanced, breakable,
    arc=5pt, outer arc=5pt,
    colback=grisPapel, colframe=grisLinea,
    boxrule=0.8pt,
    fonttitle=\bfseries\small, coltitle=azulNoche,
    top=6pt, bottom=6pt, left=10pt, right=10pt,
  },
  cajaFortaleza/.style={
    enhanced, breakable,
    arc=6pt, outer arc=6pt,
    colback=verdeSignal!8!white, colframe=verdeSignal,
    boxrule=1pt,
    borderline west={4pt}{0pt}{verdeSignal},
    fonttitle=\bfseries\small, coltitle=verdeSignal!70!black,
    drop fuzzy shadow=verdeSignal!20!white,
    top=6pt, bottom=6pt, left=10pt, right=10pt,
  },
  cajaMejora/.style={
    enhanced, breakable,
    arc=6pt, outer arc=6pt,
    colback=naranjaVivo!8!white, colframe=naranjaVivo,
    boxrule=1pt,
    borderline west={4pt}{0pt}{naranjaVivo},
    fonttitle=\bfseries\small, coltitle=naranjaVivo!80!black,
    drop fuzzy shadow=naranjaVivo!20!white,
    top=6pt, bottom=6pt, left=10pt, right=10pt,
  },
  cajaRecomendacion/.style={
    enhanced, breakable,
    arc=6pt, outer arc=6pt,
    colback=cyanNeon!6!white, colframe=cyanNeon!60!azulNoche,
    boxrule=1pt,
    borderline west={4pt}{0pt}{cyanNeon},
    fonttitle=\bfseries\small, coltitle=azulNoche,
    drop fuzzy shadow=azulNoche!20!white,
    top=6pt, bottom=6pt, left=10pt, right=10pt,
  },
}

% ── Ícono + texto (encabezados de sección y elementos) ────────────────────────
\newcommand{\iconotexto}[2]{\textcolor{cyanNeon}{\faIcon{#1}}~#2}

% ── Listados de código: estilo dark + auto-envoltura ─────────────────────────
\lstdefinestyle{estiloCodigo}{
    backgroundcolor=\color{codigoFondo},
    basicstyle=\ttfamily\small\color{white},
    commentstyle=\color{codigoComentario}\itshape,
    stringstyle=\color{codigoCadena},
    keywordstyle=\color{codigoKeyword}\bfseries,
    numberstyle=\tiny\color{codigoComentario},
    numbers=left,
    numbersep=10pt,
    showstringspaces=false,
    breaklines=true,
    breakatwhitespace=false,
    tabsize=4,
    frame=none,
    captionpos=b,
    aboveskip=6pt,
    belowskip=4pt,
    xleftmargin=14pt,
    framexleftmargin=14pt,
    literate=
      {á}{{\'a}}1 {é}{{\'e}}1 {í}{{\'i}}1 {ó}{{\'o}}1 {ú}{{\'u}}1
      {Á}{{\'A}}1 {É}{{\'E}}1 {Í}{{\'I}}1 {Ó}{{\'O}}1 {Ú}{{\'U}}1
      {ñ}{{\~n}}1 {Ñ}{{\~N}}1 {ü}{{\"u}}1 {Ü}{{\"U}}1
      {¿}{{?`}}1 {¡}{{!`}}1
}
\lstdefinestyle{estiloPython}{
    style=estiloCodigo,
    language=Python,
}
\lstset{style=estiloCodigo}
\BeforeBeginEnvironment{lstlisting}{%
  \begin{tcolorbox}[
    enhanced, breakable,
    arc=6pt, outer arc=6pt,
    colback=codigoFondo, colframe=azulNoche,
    boxrule=1pt,
    drop fuzzy shadow=azulNoche!25!white,
    top=2pt, bottom=2pt, left=0pt, right=0pt,
  ]%
}
\AfterEndEnvironment{lstlisting}{\end{tcolorbox}}

% ── Cajas de contenido temáticas ──────────────────────────────────────────────
\newtcolorbox{cajaRecuerda}{
    enhanced, breakable,
    arc=6pt, outer arc=6pt,
    colback=azulNoche!10!grisPapel,
    colframe=azulNoche, boxrule=1pt,
    borderline west={4pt}{0pt}{cyanNeon},
    fonttitle=\bfseries\small\color{white},
    title={\faIcon{info-circle}~Recuerda},
    attach boxed title to top left={yshift=-3mm, xshift=6mm},
    boxed title style={colback=azulNoche, arc=4pt, boxrule=0pt},
    drop fuzzy shadow=azulNoche!25!white,
    left=10pt, right=10pt, top=8pt, bottom=8pt,
}

\newtcolorbox{cajaConcepto}{
    enhanced, breakable,
    arc=6pt, outer arc=6pt,
    colback=verdeSignal!8!white,
    colframe=verdeSignal, boxrule=1pt,
    borderline west={4pt}{0pt}{verdeSignal},
    fonttitle=\bfseries\small\color{white},
    title={\faIcon{lightbulb}~Concepto clave},
    attach boxed title to top left={yshift=-3mm, xshift=6mm},
    boxed title style={colback=verdeSignal!80!black, arc=4pt, boxrule=0pt},
    drop fuzzy shadow=verdeSignal!20!white,
    left=10pt, right=10pt, top=8pt, bottom=8pt,
}

\newtcolorbox{cajaEjemplo}{
    enhanced, breakable,
    arc=6pt, outer arc=6pt,
    colback=cyanNeon!5!white,
    colframe=cyanNeon!70!azulNoche, boxrule=1pt,
    borderline west={4pt}{0pt}{cyanNeon},
    fonttitle=\bfseries\small\color{fondoOscuro},
    title={\faIcon{code}~Ejemplo},
    attach boxed title to top left={yshift=-3mm, xshift=6mm},
    boxed title style={colback=cyanNeon, arc=4pt, boxrule=0pt},
    drop fuzzy shadow=azulNoche!20!white,
    left=10pt, right=10pt, top=8pt, bottom=8pt,
}

\newtcolorbox{cajaEjercicio}{
    enhanced, breakable,
    arc=6pt, outer arc=6pt,
    colback=azulMedio!7!white,
    colframe=azulMedio, boxrule=1pt,
    borderline west={4pt}{0pt}{azulMedio},
    fonttitle=\bfseries\small\color{white},
    title={\faIcon{pencil-alt}~Ejercicio},
    attach boxed title to top left={yshift=-3mm, xshift=6mm},
    boxed title style={colback=azulMedio, arc=4pt, boxrule=0pt},
    drop fuzzy shadow=azulNoche!20!white,
    left=10pt, right=10pt, top=8pt, bottom=8pt,
}

\newtcolorbox{cajaReto}{
    enhanced, breakable,
    arc=6pt, outer arc=6pt,
    colback=naranjaVivo!6!white,
    colframe=naranjaVivo, boxrule=1pt,
    borderline west={4pt}{0pt}{naranjaVivo},
    fonttitle=\bfseries\small\color{white},
    title={\faIcon{fire}~Reto},
    attach boxed title to top left={yshift=-3mm, xshift=6mm},
    boxed title style={colback=naranjaVivo, arc=4pt, boxrule=0pt},
    drop fuzzy shadow=naranjaVivo!20!white,
    left=10pt, right=10pt, top=8pt, bottom=8pt,
}

% ── Indicadores de dificultad ─────────────────────────────────────────────────
\newcommand{\facil}{\textcolor{verdeSignal}{\faIcon{circle}\,\textbf{Fácil}}}
\newcommand{\intermedio}{\textcolor{naranjaVivo}{\faIcon{circle}\,\textbf{Intermedio}}}
\newcommand{\dificil}{\textcolor{rojoAlerta}{\faIcon{circle}\,\textbf{Difícil}}}

% ── Cabecera específica del manual ────────────────────────────────────────────
\fancyhead[L]{\small\color{azulNoche}\textbf{Manual --- Nuevo Proyecto}}
\fancyhead[R]{\small\color{grisTexto}\itshape Arquitectura de 3 Capas}
\renewcommand{\headrulewidth}{0.6pt}
\renewcommand{\headrule}{\color{cyanNeon}\hrule width\headwidth height 0.6pt}

% ─────────────────────────────────────────────────────────────────────────────
\begin{document}

% ══ PORTADA INFográfica ══════════════════════════════════════════════════════
\bandaTitulo{Manual para Iniciar un Nuevo Proyecto}{Arquitectura de 3 Capas: Directives $\to$ Orchestration $\to$ Execution}

\begin{center}
  \vspace{2pt}
  {\large \textbf{Autor:} Agente IA (Orquestador) \hfill \textbf{Tipo:} Guía de inicio} \\[4pt]
  {\large \textbf{Fecha:} \today} \\[6pt]
  \color{cyanNeon}\rule{\linewidth}{0.4pt}
\end{center}

\tableofcontents
\newpage

% ══ 1. INTRODUCCIÓN ═══════════════════════════════════════════════════════════
\section[Introducción]{\iconotexto{info-circle}{Introducción}}

Este manual describe los pasos para iniciar un nuevo espacio de trabajo bajo la arquitectura de 3 capas (Directives $\to$ Orchestration $\to$ Execution). El objetivo es proporcionar una estructura reutilizable que maximice la fiabilidad, separando la lógica probabilística del LLM de la ejecución determinista mediante scripts especializados.

\begin{tcolorbox}[tarjetaDato, title={Arquitectura en 3 capas}]
  \begin{tabularx}{\linewidth}{>{\bfseries\color{azulNoche}}l>{\raggedright\arraybackslash}X}
    \rowcolor{azulNoche!12}\faIcon{book}~Directives & \textit{Qué hacer} --- SOPs en YAML (Layer 1).\\
    \rowcolor{white}\faIcon{brain}~Orquestación & \textit{Decide cuándo y cuál} --- validación y enrutamiento (Layer 2).\\
    \rowcolor{white}\faIcon{terminal}~Ejecución & \textit{Cómo} --- scripts deterministas (Layer 3).\\
  \end{tabularx}
\end{tcolorbox}

% ══ 2. ESTRUCTURA DE DIRECTORIOS ══════════════════════════════════════════════
\section[Estructura de Directorios]{\iconotexto{folder-open}{Estructura de Directorios}}

Crear la siguiente estructura en la raíz del nuevo proyecto:

\begin{tcolorbox}[cajaContenido, title={Árbol de directorios}]
\begin{lstlisting}
proyecto/
|-- .agent/           # Instrucciones del sistema para el agente IA
|-- directives/       # SOPs en YAML (Layer 1)
|-- execution/        # Scripts deterministas (Layer 3)
|-- .tmp/             # Archivos temporales (regenerables)
|-- .env              # Credenciales y configuracion (no committear)
|-- opencode.json     # Configuracion del agente
|-- .gitignore        # Exclusiones de seguimiento
|-- requirements.txt  # Dependencias (se completa al definir el proyecto)
\end{lstlisting}
\end{tcolorbox}

% ══ 3. PASO 1 ═════════════════════════════════════════════════════════════════
\section[Paso 1 --- Archivos del Agente]{\iconotexto{cogs}{Paso 1: Archivos del Agente (\texttt{.agent/})}}

Copiar los archivos \texttt{FRAMEWORK.md} e \texttt{INSTRUCTIONS.md} desde la plantilla ubicada en \texttt{docs/AGENTE\_IA/} a la carpeta \texttt{.agent/} del nuevo proyecto. Estos archivos definen el marco operativo y las instrucciones del agente.

\begin{tcolorbox}[cajaContenido, title={Personalizar en \texttt{INSTRUCTIONS.md}}]
\begin{itemize}[leftmargin=*]
    \item \textbf{Sección 1:} Ajustar el tono, las restricciones de entorno y el stack tecnológico (SO, gestor de paquetes, lenguajes).
    \item \textbf{Sección 8:} Definir el formato de documentación del proyecto (Markdown, LaTeX, etc.).
    \item \textbf{Sección 10:} Completar las directrices de dominio técnico: lenguaje principal, frameworks, infraestructura, seguridad y flujo de trabajo.
\end{itemize}
\end{tcolorbox}

\begin{tcolorbox}[cajaContenido, title={Actualizar en \texttt{FRAMEWORK.md}}]
\begin{itemize}[leftmargin=*]
    \item \textbf{Available Directives:} Listar los flujos de trabajo repetibles propios del proyecto.
    \item \textbf{Setup:} Documentar los pasos de instalación específicos del proyecto.
\end{itemize}
\end{tcolorbox}

% ══ 4. PASO 2 ═════════════════════════════════════════════════════════════════
\section[Paso 2 --- Configuración del Agente]{\iconotexto{file-code}{Paso 2: Configuración del Agente (\texttt{opencode.json})}}

Crear \texttt{opencode.json} en la raíz del proyecto para que el agente cargue las instrucciones desde \texttt{.agent/}:

\begin{tcolorbox}[cajaContenido, title={\texttt{opencode.json}}]
\begin{lstlisting}
{
  "$schema": "https://opencode.ai/config.json",
  "instructions": [".agent/*.md"]
}
\end{lstlisting}
\end{tcolorbox}

% ══ 5. PASO 3 ═════════════════════════════════════════════════════════════════
\section[Paso 3 --- Primeras Directivas]{\iconotexto{clipboard-list}{Paso 3: Crear las Primeras Directivas (\texttt{directives/})}}

Cada directiva es un archivo YAML que describe un flujo de trabajo repetible. Estructura recomendada:

\begin{tcolorbox}[cajaContenido, title={\texttt{directives/mi\_flujo.yaml}}]
\begin{lstlisting}
# directives/mi_flujo.yaml
goal: "Descripcion del objetivo"
required_inputs:
  - name: "entrada_1"
    description: "Que datos se necesitan"
steps:
  - script: "execution/mi_script.py"
    mapping:
      entrada_1: "{entrada_1}"
expected_outputs: "Descripcion del resultado esperado"
edge_cases:
  - "Que hacer si algo falla"
\end{lstlisting}
\end{tcolorbox}

Crear al menos una directiva inicial para validar que el flujo funciona. Ejemplos comunes: \texttt{scrape\_website.yaml}, \texttt{research\_topic.yaml}, \texttt{git\_update.yaml}.

% ══ 6. PASO 4 ═════════════════════════════════════════════════════════════════
\section[Paso 4 --- Scripts de Ejecución]{\iconotexto{terminal}{Paso 4: Crear los Scripts de Ejecución (\texttt{execution/})}}

Cada script debe cumplir:

\begin{tcolorbox}[tarjetaIcono, title={Requisitos de un script}]
\begin{itemize}[leftmargin=*]
    \item Una sola responsabilidad.
    \item Entradas por argumentos CLI.
    \item Secretos en \texttt{.env}.
    \item Salida por stdout (preferiblemente JSON).
    \item Códigos de salida: 0 (éxito), 1+ (fallos categorizados).
    \item Validar sus propias salidas y fallar ruidosamente si algo está mal.
\end{itemize}
\end{tcolorbox}

Ejemplo mínimo en Python:

\begin{tcolorbox}[cajaContenido, title={Plantilla de script}]
\begin{lstlisting}[language=Python]
#!/usr/bin/env python3
"""Descripcion breve de la funcion del script."""
import sys

def main():
    # Validar entradas
    # Ejecutar logica
    # Validar salida
    print('{"status": "ok"}')
    return 0

if __name__ == "__main__":
    sys.exit(main())
\end{lstlisting}
\end{tcolorbox}

% ══ 7. PASO 5 ═════════════════════════════════════════════════════════════════
\section[Paso 5 --- Archivos de Configuración]{\iconotexto{wrench}{Paso 5: Archivos de Configuración}}

\subsection{.gitignore}
Crear un \texttt{.gitignore} genérico:

\begin{tcolorbox}[cajaContenido, title={\texttt{.gitignore}}]
\begin{lstlisting}
__pycache__/
*.pyc
.env
.tmp/
chroma_db/
*.aux
*.log
*.out
*.toc
venv/
.venv/
\end{lstlisting}
\end{tcolorbox}

\subsection{.env}
Plantilla para credenciales:

\begin{tcolorbox}[cajaContenido, title={\texttt{.env}}]
\begin{lstlisting}
# API Keys
# API_KEY=tu_clave_aqui
\end{lstlisting}
\end{tcolorbox}

% ══ 8. FLUJO DE TRABAJO DEL AGENTE ════════════════════════════════════════════
\section[Flujo de Trabajo del Agente]{\iconotexto{sync-alt}{Flujo de Trabajo del Agente}}

Una vez configurada la estructura, el agente opera siguiendo este ciclo:

\begin{tcolorbox}[tarjetaDato, title={Ciclo de operación}]
\begin{enumerate}[leftmargin=*]
    \item \textbf{Validación de Entorno:} Verificar que el sistema cumple los requisitos.
    \item \textbf{Búsqueda:} Revisar \texttt{directives/} para ver si existe un SOP para la tarea solicitada.
    \item \textbf{Planificación:} Seleccionar la directiva y el script de ejecución correspondiente.
    \item \textbf{Ejecución:} Invocar el script con las entradas validadas.
    \item \textbf{Validación:} Confirmar que la salida coincide con lo esperado.
    \item \textbf{Registro:} Guardar el estado en \texttt{.tmp/run\_state.json}.
    \item \textbf{Notificación:} Informar al usuario del resultado.
    \item \textbf{Aprendizaje:} Si surge un error, documentarlo y actualizar la directiva.
\end{enumerate}
\end{tcolorbox}

% ══ 9. BUENAS PRÁCTICAS ═══════════════════════════════════════════════════════
\section[Buenas Prácticas]{\iconotexto{check-circle}{Buenas Prácticas}}

\begin{tcolorbox}[cajaRecomendacion, title={Recomendaciones}]
\begin{itemize}[leftmargin=*]
    \item \textbf{Reutilizar antes de construir:} Revisar \texttt{execution/} antes de crear un nuevo script.
    \item \textbf{Directivas atómicas:} Cada directiva debe describir un solo flujo de trabajo.
    \item \textbf{Versionado de directivas:} Si cambia la lógica, conservar versiones antiguas (v1, v2).
    \item \textbf{Máximo 3 reintentos:} Si un flujo falla 3 veces, escalar al usuario.
    \item \textbf{Validar siempre:} Después de cada paso, confirmar esquema, cantidades y archivos generados.
    \item \textbf{Documentar errores:} Mantener un registro de errores comunes y sus soluciones (ej. \texttt{errores.md}) para evitar regresiones.
\end{itemize}
\end{tcolorbox}

% ══ 10. CONCLUSIÓN ════════════════════════════════════════════════════════════
\section[Conclusión]{\iconotexto{flag-checkered}{Conclusión}}

\begin{tcolorbox}[cajaFortaleza, title={Resultado}]
Con esta estructura, el nuevo proyecto queda listo para operar bajo la arquitectura de 3 capas. El esqueleto inicial es mínimo pero escalable: las directivas, scripts y configuraciones se agregan según el proyecto va tomando forma.
\end{tcolorbox}

\end{document}
